\begin{frame}{Essay 2: The Monetary Valuation of Vaccination Delay}
    \begin{itemize}
        \item \textbf{Objective}: Translating latent behavioral preferences into policy-ready \alert{Monetary Units}.
        \item \textbf{Methodology: Structural Estimation of MWTA}
        \begin{itemize}
            \item \textbf{Framework}: Random Utility Model (RUM) with a Linear-in-Incentives specification.
            \item \textbf{Estimand}: Marginal Willingness to Accept (\textbf{MWTA}) = $-\frac{\beta_{delay}}{\beta_{incentive}}$.
            \item \textbf{Identification}: Identified via experimental trade-offs between waiting time (0--6 mo) and cash (0--800 RMB).
        \end{itemize}
    \end{itemize}
    
    \begin{columns}
        \begin{column}{0.55\textwidth}
            \centering
            \includegraphics[width=\textwidth]{Subgroup differences in time sensitivity.png}
            \captionof{figure}{\footnotesize{Heterogeneity in Shadow Price (MWTA)}}
        \end{column}
        \begin{column}{0.45\textwidth}
            \begin{block}{Key Results}
                \small
                \begin{itemize}
                    \item \textbf{Average Shadow Price}: $\approx$ \textbf{111 RMB/month} ($\approx$\$15).
                    \item \textbf{The Trust Premium}: Low-trust groups require \alert{2.2x higher} compensation to tolerate delays.
                    \item \textbf{Operational ROI}: Speeding up access is more cost-effective than incentives when operational cost < MWTA.
                \end{itemize}
            \end{block}
        \end{column}
    \end{columns}
\end{frame}
